\documentclass[11pt]{article}
\usepackage[margin=1in]{geometry}
\usepackage{amsmath,amssymb,amsfonts}
\usepackage{booktabs}
\usepackage{listings}
\usepackage{xcolor}
\usepackage{hyperref}
\hypersetup{colorlinks=true,linkcolor=blue,urlcolor=blue,citecolor=blue}
\lstset{basicstyle=\footnotesize\ttfamily,frame=single,breaklines=true,columns=fullflexible}

\title{Independent Replication \#2 of \\
       WaveTrain: A Python Package for Numerical Quantum Mechanics \\
       of Chain-Like Systems Based on Tensor Trains \\
       (Riedel, Gel\ss{}, Klein, Schmidt; J.\ Chem.\ Phys.\ 158, 164801 (2023))}
\author{OpenClaw replication subagent \\
        \small session \texttt{agent:main:subagent:a1c00947-\dots}}
\date{2026-07-06 (America/Chicago)}

\begin{document}
\maketitle

\begin{abstract}
We independently re-execute the WaveTrain (Riedel et al., 2023) core Exciton-TISE
benchmark from a fresh venv, hitting the same NumPy \texttt{same\_kind} casting
bug in \texttt{scikit\_tt/solvers/evp.py} that a sibling replication also found,
and derive an equivalent minimal 4-line patch independently.  We sweep chain
length $N \in \{4,6,8,10,12\}$ and compare TT-ALS single-exciton band eigenvalues
to the closed-form tight-binding spectrum $E_k = \alpha + 2\beta\cos(2\pi k / N)$
(periodic ring), plus a bonus $N{=}6$ open-boundary test.  Core physics is
recovered to $\lesssim 2\times 10^{-3}$ Ha across the sweep, but the paper's
central TT-rank-scaling claim (C2) is only partially supported because the
\texttt{ranks{=}15} ALS ceiling saturates at $N\ge 8$.  Verdict: PARTIAL.
\end{abstract}

\section{Paper summary}
WaveTrain \cite{Riedel2023} is an open-source Python package for numerical solution
of the time-independent (TISE) and time-dependent (TDSE) Schr\"odinger equations
on quasi one-dimensional chain-like systems with nearest-neighbour (NN)
interactions.  Both Hamiltonian operators and state vectors are represented in
tensor-train (TT, or matrix-product-state) form using the authors' own
\texttt{scikit\_tt} backend.  Application classes shipped in the paper cover
pure exciton chains, pure phonon chains, coupled exciton--phonon (Holstein-like)
chains, mixed quantum--classical (Ehrenfest) dynamics, and a Bath-Map
open-system framework.  Time propagators include Strang--Marchuk (SM) splitting,
symmetrised Euler, and Krylov subspace variants; ground and excited states are
found by TT-ALS (Alternating Linear Scheme).

The paper's central algorithmic claim is that for chain-like NN Hamiltonians
the TT bond ranks of the resulting eigenstates grow only marginally with chain
length $N$---often near-linearly---so that the numerical cost of TT-based
methods grows only slightly more than linearly in $N$, while the full
Hilbert-space dimension is $2^N$ (two-level sites) and hence quickly intractable.

\section{Claims table (C1--C5)}
\begin{center}
\small
\begin{tabular}{l p{7.5cm} l c c l}
\toprule
ID & Claim & Type & Testable? & Tested? & Verdict \\
\midrule
C1 & TT-ALS recovers single-exciton band of $N{=}6$ periodic homogeneous ring
     (paper's bundled \texttt{tise\_1.py} config, $\alpha{=}0.1$, $\beta{=}-0.01$) &
     Correctness & yes (analytic $E_k$) & \checkmark & Supported \\
C2 & TT bond ranks of chain-like NN eigenstates grow only marginally with $N$
     $\Rightarrow$ wall-clock scales $\approx$ linearly in $N$ &
     Scaling & partially & \checkmark & Partial (rank cap saturates) \\
C3 & Shipped software (\texttt{wave\_train} + \texttt{scikit\_tt} from GitHub)
     is installable and runs the bundled test scripts out of the box &
     Artifact & yes & \checkmark & Fails as stated (needs 4-line patch) \\
C4 & Coupled exciton--phonon TDSE with SM propagator preserves norm and
     reproduces coherent-transport observables (Fig.\ 5--7) &
     Physics & in principle & --- & Untested (out of budget) \\
C5 & Bath-Map methodology reproduces open-system dynamics &
     Physics & nontrivial & --- & Untested (out of scope) \\
\bottomrule
\end{tabular}
\end{center}

Testable claims covered: 3 of 5 (100\% of the ones directly addressed by the
bundled Exciton test-script family).

\section{Method}

\subsection{Environment}
Fresh Python 3.12.13 virtualenv on CherryRd (macOS 25.3.0).  Dependencies:
\texttt{numpy 1.26.4} (pinned $<$2 because the paper predates NumPy 2.x and
the codebase uses removed aliases), \texttt{scipy 1.17.1}, \texttt{matplotlib}
(unused).  \texttt{wave\_train} at git HEAD, \texttt{scikit\_tt} at git HEAD
plus a 4-line dtype patch (see \S\ref{sec:patch}).

\subsection{Independent benchmark driver}
Because \texttt{TISE.eigen\_values} is only populated by the \texttt{solver='qe'}
quasi-exact code path (WaveTrain does not persist the full ALS spectrum on the
object), we monkey-patch \texttt{TISE.update\_solve} in the bench driver to
capture per-state energy (via \texttt{dyn.e\_est}) and \texttt{dyn.psi.ranks}
after every iteration.  Per-state wall-clock is recorded too.

The bench sweeps $N \in \{4,6,8,10,12\}$ at \texttt{n\_levels}${}=N+1$ (vacuum
ground state plus the $N$ single-exciton band states) with fixed
\texttt{ranks{=}15, repeats{=}20, conv\_eps{=}1e{-}8}, and compares to the
closed-form reference
\[
E_k^{\rm ref} = \eta \cup \bigl\{\alpha + 2\beta\cos(2\pi k / N) : k=0,\dots,N-1\bigr\},
\]
sorted and truncated to \texttt{n\_levels}.  Analogous open-boundary reference
$E_k = \alpha + 2\beta\cos(\pi k/(N+1))$ for $k=1,\dots,N$ is used for the
bonus open-chain test.

\subsection{The 4-line \texttt{scikit\_tt} patch}
\label{sec:patch}
The Wielandt-shift deflation in \texttt{scikit\_tt/solvers/evp.py} (line 381 of
the HEAD version we cloned) does
\begin{lstlisting}[language=Python]
micro_op += shift*tmp.dot(np.conjugate(tmp.T))
\end{lstlisting}
with \texttt{micro\_op} initialised as real but \texttt{tmp} complex, which
NumPy $\ge$1.20 refuses under \texttt{same\_kind} casting rules.  Independent
fix:
\begin{lstlisting}[language=Python]
addend = shift*tmp.dot(np.conjugate(tmp.T))
if np.iscomplexobj(addend) and not np.iscomplexobj(micro_op):
    micro_op = micro_op.astype(np.complex128, copy=False)
micro_op = micro_op + addend
\end{lstlisting}
Backup at \texttt{evp.py.bak}.  The sibling replication converged on the same
fix; this run derived it independently from the traceback.

\section{Results vs paper}

\subsection{C1: N=6 periodic ring vs analytic tight-binding}
For the paper's bundled \texttt{tise\_1.py} configuration
($N{=}6, \alpha{=}0.1, \beta{=}-0.01, \eta{=}0$) at $n_{\rm levels}=7$:
\begin{center}
\small
\begin{tabular}{l r r}
\toprule
& $E_{\rm num}$ (Ha) & $E_{\rm ref}$ (Ha) \\
\midrule
Vacuum   & $0.0000000$ & $0.0000000$ \\
Band $k{=}3$ & $0.0800000$ & $0.0800000$ \\
Band $k{=}2,4$ & $0.0900025$ & $0.0900000$ \\
Band $k{=}2,4$ & $0.0900025$ & $0.0900000$ \\
Band $k{=}1,5$ & $0.1100010$ & $0.1100000$ \\
Band $k{=}1,5$ & $0.1100460$ & $0.1100000$ \\
Band $k{=}0$ & $0.1199000$ & $0.1200000$ \\
\bottomrule
\end{tabular}
\end{center}
Max error $9.99 \times 10^{-5}$ Ha, RMS $4.38 \times 10^{-5}$ Ha.  This matches
WaveTrain's own internally-printed ``Difference'' field (which is the paper's
implicit accuracy target).  Claim \textbf{C1 supported}.

\subsection{C2: scaling sweep}
\begin{center}
\small
\begin{tabular}{r r r r r}
\toprule
$N$ & wall\_s & max\_err (Ha) & RMS (Ha) & max\_bond\_rank \\
\midrule
4  & 0.3    & $6.31\times 10^{-15}$ & $3.66\times 10^{-15}$ & 4  \\
6  & 3.0    & $9.99\times 10^{-5}$  & $4.38\times 10^{-5}$  & 8  \\
8  & 136.6  & $1.92\times 10^{-3}$  & $1.02\times 10^{-3}$  & 15 (saturated) \\
10 & 112.1  & $9.67\times 10^{-15}$ & $3.35\times 10^{-15}$ & 15 (saturated) \\
12 & 1026.4 & $2.77\times 10^{-4}$  & $1.08\times 10^{-4}$  & 15 (saturated) \\
\midrule
\multicolumn{5}{l}{Bonus: $N{=}6$ open chain: wall 3.6 s, max\_err $1.02\times 10^{-3}$.} \\
\bottomrule
\end{tabular}
\end{center}

Canonical numerical values in \texttt{report/evidence/bench.json}.

Observations:
\begin{enumerate}
\item Accuracy is \emph{non-monotone} in $N$ (N=4 and N=10 hit machine
      precision; N=6 and N=8 do not).  Hypothesis: n\_levels${}=N{+}1$ fits
      cleanly inside the single-exciton band exactly when $N+1 \le N+1$
      trivially, but deflation for higher states starts pushing the ALS
      iterate toward two-exciton subspaces whose intrinsic TT rank exceeds
      the \texttt{ranks{=}15} ceiling.  Discussion in
      \texttt{report/failure\_analysis.md} \S3 and open question Q2.
\item Wall-clock grows sharply for the last few ALS eigenstates at each $N$
      (the deflation stack grows with each found state).  At $N{=}12$ the last
      three eigenstates each cost $\sim 150$--$200$ s.
\item TT bond-rank profile at $N{=}12$ is
      $[1,2,4,8,15,15,15,15,15,8,4,2,1]$: the middle bonds are saturated at
      the user-set ceiling.  This means we can conclude
      ``$r{=}15$ suffices for the single-exciton band across $N \in [4,12]$
      to accuracy $\lesssim 2\times 10^{-3}$ Ha'' but not the stronger
      ``TT ranks themselves grow only marginally with $N$'' claim.  Claim
      \textbf{C2 partial}.
\end{enumerate}

\subsection{C3: software artifact}
Bundled \texttt{tise\_1.py} crashes as-shipped on NumPy $\ge$1.20 with the
casting bug detailed in \S\ref{sec:patch}.  After a 4-line patch, it runs to
completion and reproduces the paper's expected accuracy.  Claim
\textbf{C3 fails as-stated} (needs the patch).  Note: the sibling replication
had already documented this, and we independently confirm.

\subsection{Bonus: open-chain}
$N{=}6$ open chain at $n_{\rm levels}{=}7$: max error $1.02\times 10^{-3}$ Ha,
RMS $4.56\times 10^{-4}$ Ha vs the open-boundary analytic reference
$E_k = \alpha + 2\beta\cos(\pi k/(N+1))$.  Serves as a sanity check that the
analytic reference (and the underlying tight-binding physics) switches
correctly under boundary-condition change.  Slightly worse accuracy than
periodic $N{=}6$ (1e-3 vs 1e-4) is expected because open-boundary
eigenstates have larger site-dependence and hence higher TT bond ranks.

\section{Verdict}
LLM-judge output (GPT-5.4 via localhost:44497 Argo proxy, free endpoint; Opus
4.7/4.8 broken by upstream response-parsing bug on 2026-07-06, fallback to
GPT-5.4 which is equally free-tier). See \texttt{report/evidence/llm\_judge.json}
for the full rubric.

\textbf{Verdict: PARTIAL.}  Coverage 100\%, agreement 67\%, confidence high.

Judge's per-claim scoring:
\begin{itemize}
\item \textbf{C1 (partial):} Excellent at $N{=}4$ and $N{=}10$, good at $N{=}6$
      ($\sim$1e-4), noticeably worse at $N{=}8$ (1.9e-3) and $N{=}12$ (2.8e-4).
      Recovery of the single-exciton band is qualitative but not uniformly
      ``high accuracy''.
\item \textbf{C2 (partial):} Observed TT ranks capped by the ALS ceiling and
      saturated at $N \ge 8$: shows $r{=}15$ suffices, not that intrinsic
      ranks grow only marginally with $N$. Wall-clock rises from 0.3\,s
      (N=4) to 1026\,s (N=12) with severe late-deflation slowdown---far from
      near-linear over the tested range, even though vastly better than
      brute-force diagonalisation in principle.
\item \textbf{C3 (partial):} Bundled software does not run on Python 3.12 /
      NumPy 1.26 as-shipped; requires 4-line patch. A real reproducibility
      defect even though the fix is minimal.
\end{itemize}

C4 and C5 untested (TDSE and Bath-Map out of budget for a one-shot
replication).

\section{Open Questions}
The five heavy-duty open questions this replication surfaces, with concrete
next steps, are stored as machine-readable JSON at
\texttt{report/open\_questions.json}.  Summaries:
\begin{description}
\item[Q1.] Do TT bond ranks of single-exciton eigenstates \emph{actually} grow
      linearly in $N$, or is $r{=}15$ already a plateau?  Ranks-vs-accuracy
      elbow study needed.
\item[Q2.] Why is $N{=}8$ max\_err worse than $N{=}6$ and $N{=}10$?
      Deflation into two-exciton subspaces beyond $r{=}15$ is the leading
      hypothesis.
\item[Q3.] Is the NumPy \texttt{same\_kind} bug historically silent (older
      NumPy accepted the cast by discarding imaginary parts)?  If so, how much
      accuracy did \texttt{scikit\_tt} users silently lose pre-2020?
\item[Q4.] Does the exciton--phonon TDSE path also require this patch?
      (Same \texttt{evp.py} code path may be reached from the
      \texttt{Exc\_Pho\_Coupling} scripts.)
\item[Q5.] Cross-check TT-ALS at $N{=}12$ against a direct
      \texttt{scipy.sparse.linalg.eigsh} on the assembled 4096$\times$4096
      Hamiltonian.
\end{description}

\begin{thebibliography}{1}
\bibitem{Riedel2023}
J.~Riedel, P.~Gel\ss{}, R.~Klein, B.~Schmidt,
\newblock ``WaveTrain: A Python package for numerical quantum mechanics of
chain-like systems based on tensor trains,''
\newblock \emph{J.\ Chem.\ Phys.} \textbf{158}, 164801 (2023).
\newblock DOI: \href{https://doi.org/10.1063/5.0147314}{10.1063/5.0147314}.
\newblock arXiv preprint \href{https://arxiv.org/abs/2302.03725}{2302.03725}.
\end{thebibliography}

\end{document}
